\documentclass[10pt]{article}                         
\pagestyle{plain} %-- Document Set Up Command

%-- Document Setup & Layout
\usepackage{fancyhdr}
\usepackage{comment}
\usepackage{multicol}
\usepackage{geometry}
\usepackage{setspace}


%-- Header/Footer
\pagestyle{fancy}
\fancyhf{} % Clear all header and footer fields
\fancyhead[L]{Jack Lingbeck} % Left header with name
\fancyhead[R]{April 18th, 2026} % Right header with date

\begin{document}

\begin{center}
    \Large \textbf{Math AutoBiography} \\
    \vspace{0.2cm}
\end{center}

Instructions. Write a reflection of approximately 300–600 words (about 1-2 pages double-spaced), or record an
audio/video response of 2 - 4 minutes. Address the prompts below by choosing at least three.
Suggested Prompts (choose a minimum of three from this list):
\begin{enumerate}
\item Describe your current relationship with mathematics. How do you feel about math now, and what experiences
have shaped that feeling?
\item Reflect on a turning point in your math learning – an experience that significantly increased or decreased your
confidence.
\item What messages have you received (from teachers, family, peers, society) about who is good at math? How have
those messages influenced your math identity?
\item When math feels difficult, what strategies or supports help you? What has not helped? (e.g., study groups,
practice types, explanations, office hours)
\item What helps you feel comfortable asking questions or making mistakes? Conversely, what makes it harder?
\item What are your goals for this course? What would you like to be different about your experience with math by the
end of the semester?
\item Share an example of how you use mathematical thinking outside school, even if you don’t label it as “math."
\end{enumerate}

\noindent
\begin{doublespacing}
\noindent
1/2) Mathmatics is a subject that has always been a part of my life, but my relationship with it has been complex. 
When I was a kid, math was kind of boring and felt the same thing over and over again, until I hit Algebra 1.
Oh hey new letters and formulas that I have to learn, this is different. I was actually interested in learning it
and I did well in it, this year was a pivotal moment for me. It made me feel like I could do math and that it was something 
I could understand and even play around with. Then came algebra 2 and trigonometry, it didn't have that same magic and
I was struggling to keep up or keep my interest. Years later I enrolled in community college and started Calculus 1.
Damn, this damn brought back the magic tenfold. I was actually thrilled to learn about limits, derivatives, and integrals. 
It was like a whole new world of math opened up to me, and it was also my first 100\% in a class. Calculus 1 was 
definitely a permanent inflection point for me in my math journey, I couldn't stop loving math from that point on. 
All 3 calculus professors were wonderful, I felt like they really cared about the students and wanted us to succeed. 
I even had nice conversations with them before class started and they were constantly passionate.
\\\\
I know it may sound repetitive, but once i started East Bay, with upper division math, oh boy did
that enjoyment curve become exponential, with some dips here and there. Community college felt like all the
classes and professors were seperate, not even remotely like that for East Bay. It's a fantastic and 
welcoming community of students and professors, and it's a great place to learn and grow as a mathematician.
Because I've taken so many classes here, ive had repeat professors and I get to learn about them as people.
What im getting at, is that Math can be just a subject, the set of mathmatics I've learned and thats it,
or, how I feel now, is that Math is more about exploring, bringing people along with you, and sharing what
we all know and learn while, of course, still actually trying to pass classes. But It's impossible to
seperate Math, the friends I've made, and the amazing professors here. It's all one big squiggly domain that
can't be split up into finite subcovers, I need every detail, every experience, and every helpful person to
build what math truely means to me. It's much more difficult to struggle when everyone wants to help and be together.
\\\\
4) Oh that feeling of math being difficult. Over the years at East Bay I've really felt the difference
between studying alone and trying to do it all alone vs handing out with friends and having study
routines and support systems. A couple of friends helped me through tough analysis, but then when the time came
i was able to help others in other classes when i felt smart and they wanted help. Even office hours help
way more than i always expect, I barely attended any during community college or my undergrad,
but now when i go to them, the professors are always welcoming and ready to answer anything or to clarify things.
It's hard to say what has not helped, since all of these things can easily be twisted or misused or not used
to their fullest potential. For example, study groups can be a huge help, office hours can be a huge help, 
but if you don't have people you mesh well with, or if a professor is busy with others or having such fruitful 
conversations with the obviously smarter students, it can deter someone, even me at times, from using these resources.

\vspace{6cm}
5) It honestly took years for me to be comfortable asking questions in math class. 
I think it was because I had this fear of looking stupid or looking too smart or being the only one who 
didn't understand something, as i figure most people do at some point. Sometimes i felt like i didnt want to ask 
a niche or specific question because i thought it would be too much for the teacher to answer or that it would 
take up too much time. What changed was that, because i always sat up front, i didn't realize much others
in the back of the class could be struggling, so if i asked a question, maybe it could help others.
Or hell, I could even ask a question i fully knew the answer to just so the rest of the class had some
helpful info or revelation that they might not have gotten. 
\\\\
Honestly, having such amazing professors in my whole college time has helped me be comfortable asking questions. 
They were always so passionate and helpful, and they made it clear that they wanted us to succeed.
The way they handle any question, from potentially dumb, to regular, to super specific, 
they were always happy to anwer. I carried this mentality with me when I was a TA, just answer questions and
the students will appriciate it, and even point out how it was a good question. I think what makes it harder
is a few lower division general ed professors wouldn't have that...for lack of better words...crack head energy.
To be so enveloped in their field that you can feel that they want to teach mostly to share their knowledge
and show off how damn cool math can be. If i feel that a professor doesn't want to have that connection
with their students, i care less to ask questions and when i really want to, i worry about their response .

\end{doublespacing}
\vspace{1cm}
- Jack Lingbeck \\\\\\
\noindent
PS: Man I wish Latex had a spell checker, appologies for random basic typos.
\\\\\\
PPS: Reached a good flow state and probably typed too much, but this was nice


\end{document}